\documentclass[11pt]{article}
\usepackage[margin=0.95in]{geometry}
\usepackage[T1]{fontenc}
\usepackage[utf8]{inputenc}
\usepackage{lmodern}
\usepackage{microtype}
\usepackage{amsmath,amssymb}
\usepackage{booktabs}
\usepackage{tabularx}
\usepackage{array}
\usepackage{xcolor}
\usepackage{enumitem}
\usepackage[hidelinks]{hyperref}
\usepackage[capitalize,noabbrev]{cleveref}
\usepackage{caption}

\definecolor{obsblue}{HTML}{075FA8}
\definecolor{badred}{HTML}{B42318}
\newcommand{\obs}[1]{\textcolor{obsblue}{#1}}
\newcommand{\bad}[1]{\textcolor{badred}{#1}}
\setlist[itemize]{leftmargin=*,topsep=2pt,itemsep=2pt}
\setlist[enumerate]{leftmargin=*,topsep=2pt,itemsep=2pt}
\captionsetup{font=small,labelfont=bf}

\title{\textbf{Pair the target with its own noise}\\[3pt]
\large What actually helps when a few-step generator copies selected teacher images}
\author{}
\date{}

\begin{document}
\maketitle
\vspace{-0.25in}

\begin{abstract}
\noindent
A slow image generator makes four images for one caption. We train a fast generator to copy them.
This report asks what makes that work, and finds that most of the answer depends on the training
budget and on how often a caption is seen again. Measuring one factor at a time, three seeds each:
pairing each chosen image with \textbf{the noise that produced it} is worth \obs{$+0.083$} on
T2I-CompBench++ after \obs{6{,}000} updates, but only \obs{$+0.031$} after \obs{40{,}000}, where a
plain label does as well. Keeping the chosen image \textbf{the same} between visits is worth
\obs{$+0.042$} at \obs{3{,}000} captions and \bad{nothing} at \obs{24{,}000}, because it can only act
when a caption is revisited. \textbf{Which} image is chosen is worth at most \obs{$+0.037$}. The
order in which several images are visited is worth nothing at all. Two results survive scaling.
Fixing the target buys alignment by \textbf{collapsing} what a caption can produce, failing a
pre-registered diversity gate at \bad{$-0.167$}, while coupling raises alignment and variation
together. And reward selection at \textbf{sampling} time --- best of four, no training --- beats the
same signal used as a training label. The coupling method is not new; it is text-conditioned Reflow.
The contribution is the diagnosis and the scope.
\end{abstract}


\section{Summary}

\Cref{sec:goal} states the question we set out to answer and \cref{sec:story} tells what happened to
it. \Cref{tab:summary} lists every factor we varied, each measured with the others held still, at
the budget and pool size the campaign used throughout. \textbf{Read it with \cref{sec:scope}}, which
shows that most of these numbers move --- some of them to zero --- when the budget or the number of
captions changes.


\begin{table}[h]
\centering
\caption{Every factor, measured with the others held still. Paired bootstrap over seeds and test
captions. Three seeds per arm.}
\label{tab:summary}
\small
\begin{tabular}{@{}llrl@{}}
\toprule
Factor & Change & Effect & 95\% interval \\
\midrule
Training recipe & flow matching $\to$ consistency & \obs{$+0.2043$} & \obs{$[+0.1770,+0.2254]$} \\
\textbf{Coupling} & fresh noise $\to$ the target's own noise, one target & \obs{$\mathbf{+0.0832}$} & \obs{$[+0.0608,+0.1049]$} \\
\textbf{Coupling} & fresh noise $\to$ each target's own noise, four targets & \obs{$\mathbf{+0.0863}$} & \obs{$[+0.0546,+0.1194]$} \\
Persistence & target stops changing, \obs{40}k steps & \obs{$+0.0954$} & \obs{$[+0.0797,+0.1097]$} \\
Persistence & target stops changing, \obs{6}k steps & \obs{$+0.0417$} & \obs{$[+0.012,+0.072]^{\ast}$} \\
Label & random $\to$ VQAScore top-1, \obs{40}k steps & \obs{$+0.0366$} & \obs{$[+0.0268,+0.0461]$} \\
Targets & four $\to$ one, churn held at \obs{100\%} & \obs{$+0.0266$} & \bad{$[-0.0028,+0.0580]$} \\
Targets & four $\to$ two, churn held at \obs{100\%} & \obs{$+0.0275$} & \bad{$[-0.0028,+0.0594]$} \\
Order & churn \obs{14\%} $\to$ \obs{100\%}, two targets & \bad{$+0.0017$} & \bad{$[-0.0176,+0.0212]$} \\
Order & churn \obs{43\%} $\to$ \obs{100\%}, four targets & \bad{$-0.0019$} & \bad{$[-0.0420,+0.0408]$} \\
\bottomrule
\end{tabular}

\vspace{2pt}
{\footnotesize $^{\ast}$widened for the caption-to-target map, which a seed-only bootstrap does not
resample (\cref{sec:assign}). Every other interval here is seed-and-prompt only.}
\end{table}


\section{What we set out to do}
\label{sec:goal}

A slow generator makes several images for one caption. They are not equally good: one may show both
objects the caption asks for, another may miss one. A fast generator is then trained to copy the
slow one. The obvious idea is to copy \emph{the best} image rather than a random one, and to get a
better fast generator for free.

To do that you must know which image is best. The reliable way is to decode every image and ask a
vision--language model, which is slow. So we built a cheap score, \textbf{PRS}, that ranks images
\emph{inside} the generator with no decode step (\cref{sec:scores}), and set out to answer:

\begin{enumerate}[nosep]
  \item Does copying a better image make a better fast generator?
  \item Can a cheap score find the better image well enough to matter?
\end{enumerate}

Both answers turned out to be ``barely'', and the experiments that established that led somewhere
more useful. \Cref{sec:story} tells it in order, because the final result only makes sense against
what it replaced.

\section{How the answer changed, in four stages}
\label{sec:story}

\paragraph{Stage 1. Choosing barely matters, but \emph{committing} to a choice helps a lot.}
We compared rules that pick one of four teacher images per caption. Making the label better --- from
a random pick to the VQAScore top-1 of four --- was worth almost nothing at \obs{6{,}000} training steps,
\obs{$+0.0042$} with an interval crossing zero. But a rule that picked \emph{at random} and then
\emph{kept its choice} beat a rule that redrew on every visit by \obs{$+0.0417$}
(\cref{sec:stage1}). The gain had nothing to do with picking well. We called this
\textbf{persistence} and assumed the mechanism was stability: a target that stops moving is easier
to regress onto.

\paragraph{Stage 2. That name was wrong. Order does not matter at all.}
If stability were the mechanism, how \emph{often} the target changes should matter. We fixed the
exact multiset of images a caption trains on and varied only the order, from changing once to
changing every single visit. We did not detect any effect, with the switch rate varied sevenfold
(\cref{sec:stage2}). A gradient-variance measurement agreed --- switching accounts for only
\obs{11.2\%} of the variance, too little to explain the effect. So the benefit of ``fixing the
target'' was not about the target being fixed.

\paragraph{Stage 3. The real cause is which \emph{noise} is paired with the image.}
Every teacher image was produced from a specific noise sample: the teacher started at
$z_0$ and arrived at $x_0$. Our training loss discarded that pairing and re-noised the chosen image
with a fresh random draw, asking the student to learn a straight line between two points the teacher
never connected. Restoring the pairing is worth \obs{$+0.0832$}, twice the persistence effect and
the largest thing here apart from the choice of objective (\cref{sec:stage3}).

\paragraph{Why the earlier result looked like persistence.}
Every arm that ``fixed the target'' also happened to use fewer distinct images, and fewer distinct
images means fewer distinct \emph{mismatched} noise--image pairs. Fixing the target was an indirect,
partial way of reducing the harm done by broken coupling. Fixing the coupling directly is better,
and it removes the reason to discard images at all: \textbf{all four teacher images, each paired with
its own noise, beat one image with random noise by \obs{$+0.0597$}}. So the original goal --- keep
the best image, throw the rest away --- was solving the wrong problem.

\paragraph{Stage 4. And then scale moved almost all of it.}
Every number above is from \obs{6{,}000} updates over \obs{3{,}000} captions. Run the same arms
longer, or over more captions, and the picture changes: the persistence gain vanishes when captions
stop being revisited, the coupling gain more than halves, and the label gain grows
(\cref{sec:scope}). Two findings survive. Fixing the target buys its alignment by collapsing what a
caption can produce, while coupling does not (\cref{sec:diversity}); and reward selection applied at
sampling time beats the same signal used as a training label (\cref{sec:refs}).

\paragraph{And the cheap score does not earn its place.}
Trained downstream, PRS is beaten by CLIP-I and is indistinguishable from using no score at all
(\cref{sec:scores}). Skipping the decode also saves less than it sounds, because the decoder costs
less than the two extra network calls the score needs.

\section{Setup}

\subsection{Model and data}
Stable Diffusion 3.5-Medium. A \textbf{VAE} turns a $512\times512$ image into a
$16\times64\times64$ \emph{latent} and back. Three \textbf{text encoders} turn the caption into
vectors. A \textbf{denoiser} of 24 transformer blocks and 2.24 billion weights predicts how to move
a latent. Only the denoiser is trained.

Training uses \obs{3{,}000} COCO \texttt{train2017} captions, each with the photograph it was
written for. We removed \obs{4{,}190} captions that repeat inside the split and \obs{179} that also
appear in \texttt{val2017}. For each caption the frozen slow generator makes $N=\obs{4}$ images at
\obs{8} steps and guidance \obs{7.0}. Image $j$ of caption $i$ starts from noise
$z_0^{(i,j)}$ drawn with seed $1000\,i+j$; the trainer uses the same rule, so it always rebuilds the
image that was scored.

\subsection{The flow}
The denoiser learns a straight path between a latent $x_0$ and noise $\varepsilon$:
\begin{equation}
  x_\sigma = (1-\sigma)x_0 + \sigma\varepsilon, \qquad
  \frac{\mathrm{d}x_\sigma}{\mathrm{d}\sigma} = \varepsilon - x_0 .
\label{eq:path}
\end{equation}
It predicts that speed as $v_\theta(x_\sigma,\tau,c)$, where $c$ is the caption and $\tau=T\sigma$
with $T=\obs{1000}$, the \texttt{num\_train\_timesteps} of the released model. $T$ is a property of
the checkpoint, not a free choice. Images are made with Euler steps,
$x_{k+1}=x_k+(\sigma_{k+1}-\sigma_k)v_\theta$. The slow generator runs the denoiser twice per step
and mixes the results, $g=v^-+w(v^+-v^-)$ with $w=\obs{7}$; our fast generator runs one call per
step and learns the mixed result directly.

\subsection{Two objectives}
\label{sec:losses}
Which objective runs is set by the arm name. An arm from one group must never be compared with an
arm from the other.

\paragraph{FM, endpoint flow matching.}
Take the chosen image's endpoint latent $x_0$, add fresh noise at a random level, ask for the
direction back out:
\begin{equation}
  u\sim\mathcal U(0,1),\quad \sigma=\frac{su}{1+(s-1)u},\ s=\obs{3.0},\qquad
  \boxed{\;\mathcal L_{\mathrm{FM}}=\bigl\|v_\theta\bigl((1-\sigma)x_0+\sigma\varepsilon,\,T\sigma,\,c\bigr)-(\varepsilon-x_0)\bigr\|_2^2\;}
\label{eq:lossfm}
\end{equation}
averaged over \obs{5} draws per update. \textbf{The choice of $\varepsilon$ is the subject of
\cref{sec:stage3}.}

\paragraph{CD, consistency distillation.}
Here the teacher's whole path is used. With states $z_0..z_8$, teacher indices
$\mathcal W=\{3..7\}$ and a gap $\Delta\in\{1,2,3\}$, the student at the noisier state $k_s=k_t-\Delta$
predicts $\hat x_0=z_{k_s}-\sigma_{k_s}v_\theta$, against the teacher's
$\tilde x_0=z_{k_t}-\sigma_{k_t}(z_{k_t+1}-z_{k_t})/(\sigma_{k_t+1}-\sigma_{k_t})$, under a
pseudo-Huber distance with $c=\obs{0.13824}$.

\subsection{Training and testing}
\obs{6{,}000} updates unless stated, AdamW at \obs{$10^{-5}$}, \obs{300} warm-up, gradient clipping
\obs{1.0}, \obs{5} supervised pairs per update, $512\times512$. Every arm shares prompts, order,
optimiser, schedule, initialisation and compute; only the target and its noise differ.

Testing uses a \textbf{one-sample variant} of T2I-CompBench++
(\texttt{Karine-Huang/T2I-CompBench} at \texttt{1b70949}), \obs{2{,}398} prompts, the official
per-category scorers, headline the macro average of the 8 categories. \textbf{The official protocol
generates ten images per prompt; we generate one.} Our numbers are therefore internally comparable
across arms, which is what this report needs, but they are \emph{not} comparable with published
T2I-CompBench++ scores and should not be quoted as such. \textbf{GenEval2} (\texttt{facebookresearch/GenEval2} at \texttt{a6e82d2}), \obs{800}
prompts as shipped, Soft-TIFA. \textbf{One} image per prompt, \obs{4} steps, guidance \obs{1.0},
shift \obs{3.0}. Intervals are a paired bootstrap over seeds and captions, captions resampled within
category.


\section{Reference points}
\label{sec:refs}

Without these rows a student's score cannot be placed against anything.

\begin{table}[h]
\centering
\caption{What the models we distil from actually score, at the same test settings.}
\label{tab:refs}
\small
\begin{tabular}{@{}llrrr@{}}
\toprule
Model & Sampling & Calls & CompBench++ & GenEval2 \\
\midrule
SD3.5-M, best of \obs{8} & \obs{28} steps, guidance \obs{7} & \obs{448} & \obs{0.5495} & \obs{0.2537} \\
SD3.5-M, best of \obs{4} & \obs{28} steps, guidance \obs{7} & \obs{224} & \obs{0.5428} & \obs{0.2304} \\
SD3.5-M, one sample & \obs{28} steps, guidance \obs{7} & \obs{56} & \obs{0.5053} & \obs{0.1705} \\
Teacher (SD3.5-M) & \obs{8} steps, guidance \obs{7} & \obs{16} & \obs{0.4544} & \obs{0.1758} \\
$G$, guidance absorbed & \obs{8} steps, guidance \obs{1} & \obs{8} & \obs{0.4593} & \obs{0.1823} \\
$G$ run at 4 steps & \obs{4} steps, guidance \obs{1} & \obs{4} & \obs{0.2859} & \obs{0.1146} \\
Untrained & \obs{4} steps, guidance \obs{7} & \obs{8} & \obs{0.2557} & \obs{0.0732} \\
Untrained & \obs{4} steps, guidance \obs{1} & \obs{4} & \obs{0.1743} & \obs{0.0335} \\
\bottomrule
\end{tabular}
\end{table}

Four things follow. $G$ \textbf{matches the teacher at half the calls} (\cref{sec:g8}). Running $G$
at \obs{4} steps with no step distillation already gives \obs{0.2859}, so that --- not the untrained
model --- is the bar a 4-step student must clear. Every consistency-distilled student in this report
\emph{beats} the \obs{8}-step teacher it copied (\cref{tab:main}), so that teacher is not a ceiling.

And selection is worth more at \emph{sampling} time than anything we moved into the weights. Picking
the best of \obs{4} samples by VQAScore adds \obs{$+0.0375$} on CompBench++ and \obs{$+0.0599$} on
GenEval2, with no training at all; the same signal used as a training label added \obs{$+0.0336$}
and \obs{$+0.0366$} after \obs{40{,}000} updates. The gain is largest exactly where every trained arm
failed: \texttt{spatial} rises from \obs{0.2716} to \obs{0.3696} under best-of-\obs{4}. An earlier
campaign in this project measured best-of-\obs{4} on the \obs{8}-step teacher at \obs{$+0.0529$}
$[+0.0434,+0.0625]$; the smaller figure here is expected, since better base samples leave less to
gain. \textbf{After two campaigns, the evidence is that reward selection works at inference and does
not transfer into the weights.}


\section{Stage 1: choosing barely matters, committing helps}
\label{sec:stage1}

\begin{table}[h]
\centering
\caption{The main campaign. Three seeds unless marked. \texttt{FM} blocks are comparable with each
other; \texttt{CD} uses the other objective and is separate.}
\label{tab:main}
\small
\begin{tabular}{@{}llrrr@{}}
\toprule
Arm & Target & $\rho$ & CompBench++ & GenEval2 \\
\midrule
\multicolumn{5}{@{}l@{}}{\emph{Flow matching, \obs{6{,}000} steps}} \\
\texttt{PRS-top1}   & fixed  & \obs{0.379} & \obs{0.3022} & \obs{0.1384} \\
\texttt{VQA-fix}    & fixed  & \obs{0.685} & \obs{0.3021} & \obs{0.1299} \\
\texttt{PRS-fix}    & fixed  & \obs{0.260} & \obs{0.3001} & \obs{0.1363} \\
\texttt{VQA-top1}   & fixed  & \obs{1.000} & \obs{0.2899} & \obs{0.1324} \\
\texttt{RAND-fix}   & fixed  & \obs{0.000} & \obs{0.2857} & \obs{0.1246} \\
\texttt{PRS-soft}   & moves  & \obs{0.260} & \obs{0.2709} & \obs{0.1056} \\
\texttt{VQA-soft}   & moves  & \obs{0.685} & \obs{0.2475} & \obs{0.0958} \\
\texttt{RAND-soft}  & moves  & \obs{0.001} & \obs{0.2440} & \obs{0.0922} \\
\midrule
\multicolumn{5}{@{}l@{}}{\emph{Flow matching, \obs{40{,}000} steps}} \\
\texttt{VQA-top1}   & fixed  & \obs{1.000} & \obs{0.4131} & \obs{0.2349} \\
\texttt{PRS-top1}   & fixed  & \obs{0.379} & \obs{0.4101} & \obs{0.2189} \\
\texttt{RAND-fix}   & fixed  & \obs{0.000} & \obs{0.3765} & \obs{0.2049} \\
\texttt{VQA-soft}   & moves  & \obs{0.685} & \obs{0.3199}$^{\dagger}$ & \obs{0.1512} \\
\texttt{RAND-soft}  & moves  & \obs{0.001} & \obs{0.2810} & \obs{0.1172} \\
\midrule
\multicolumn{5}{@{}l@{}}{\emph{Consistency distillation, \obs{6{,}000} steps}} \\
\texttt{CD-VQA-top1} & fixed & \obs{1.000} & \obs{0.4943} & \obs{0.2408} \\
\texttt{CD-PRS-top1} & fixed & \obs{0.379} & \obs{0.4901} & \obs{0.2222} \\
\texttt{CD-RAND-fix} & fixed & \obs{0.000} & \obs{0.4864} & \obs{0.2231} \\
\bottomrule
\end{tabular}

\vspace{2pt}
{\footnotesize $^{\dagger}$one seed.}
\end{table}

At \obs{40{,}000} steps the VQAScore top-1 is worth \obs{$+0.0366$} $[+0.0268,+0.0461]$ over no label,
and a mediocre one (\obs{$\rho=0.379$}) already reaches \obs{$+0.0336$} $[+0.0239,+0.0431]$ of it.
\textbf{Rank quality saturates early.} At \obs{6{,}000} steps neither is separable from zero. Under
the consistency objective the VQAScore top-1 is worth \obs{$+0.0079$} $[+0.0010,+0.0149]$.


\section{Stage 2: the order of targets does not matter}
\label{sec:stage2}

To separate \emph{how often the target changes} from \emph{how many targets there are}, we fixed the
multiset of targets a caption sees and varied only the order. Each caption is visited exactly
\obs{8} times (\obs{6{,}000} updates $\times$ \obs{4} workers $\div$ \obs{3{,}000} captions), so at
support size $m$ every one of $m$ targets is used $8/m$ times in every condition. No label is used
anywhere in this section.

\begin{table}[h]
\centering
\caption{Same targets, same counts, different order. Churn is the measured share of revisits where
the target changed, averaged over \obs{3{,}000} captions.}
\label{tab:order}
\small
\begin{tabular}{@{}lrrrr@{}}
\toprule
Targets $m$ & Order & Churn & CompBench++ & GenEval2 \\
\midrule
1 & fixed        & \obs{0\%}     & \obs{0.2696} & \obs{0.1224} \\
2 & long blocks  & \obs{14.3\%}  & \obs{0.2687} & \obs{0.1217} \\
2 & shuffled     & \obs{57.1\%}  & \obs{0.2550} & \obs{0.1024} \\
2 & alternating  & \obs{100\%}   & \obs{0.2704} & \obs{0.1112} \\
4 & long blocks  & \obs{42.9\%}  & \obs{0.2449} & \obs{0.0867} \\
4 & shuffled     & \obs{85.5\%}  & \obs{0.2494} & \obs{0.0916} \\
4 & alternating  & \obs{100\%}   & \obs{0.2429} & \obs{0.0857} \\
\bottomrule
\end{tabular}
\end{table}

\begin{itemize}
  \item \textbf{Order: nothing.} At two targets, \obs{14\%} against \obs{100\%} churn gives
    \bad{$-0.0017$} $[-0.0209,+0.0180]$. At four targets, \obs{43\%} against \obs{100\%} gives
    \bad{$+0.0019$} $[-0.0413,+0.0422]$. We did not detect an order effect at either support size. The four-target interval still admits an effect as large as \obs{0.042}, so this is a failure to detect, not a demonstration that order is irrelevant.
  \item \textbf{Number of targets: marginal.} Two targets against four, both at \obs{100\%} churn,
    gives \obs{$+0.0275$} $[-0.0029,+0.0589]$, an interval that touches zero.
\end{itemize}

\paragraph{So ``persistence'' was the wrong name.}
The gain we first measured by fixing the target is not a gain from \emph{stability}. Making the
target change on every single visit costs nothing, as long as the number of distinct targets is held
down. And \cref{sec:stage3} shows even that is unnecessary when the noise is coupled.

\paragraph{A limit of this design.}
With the multiset fixed at \obs{8} visits, more targets means fewer repeats of each: $m=1$ shows one
target \obs{8} times, $m=4$ shows each twice. So ``more targets'' and ``fewer repeats per target''
cannot be separated here. Separating them needs a different budget, for example $m=4$ over
\obs{16} visits against $m=2$ over \obs{8}.


\section{Stage 3: pair each target with its own noise}
\label{sec:stage3}

\Cref{eq:lossfm} needs a noise $\varepsilon$. The obvious choice, and the one we used at first, is a
fresh draw. But every $x_0$ came from a specific noise: the teacher rolled $z_0^{(i,j)}$ forward and
arrived at $x_0^{(i,j)}$. Drawing a fresh $\varepsilon$ throws that pairing away and asks the student
to learn a straight line between an unrelated pair of points.

\begin{equation}
  \text{fresh:}\quad \varepsilon\sim\mathcal N(0,I)
  \qquad\text{against}\qquad
  \text{coupled:}\quad \varepsilon = z_0^{(i,j)} .
\label{eq:coupling}
\end{equation}

With coupling, $(1-\sigma)x_0+\sigma\varepsilon$ is the \emph{straight interpolation induced by the
teacher's source--endpoint pairing}. It is not the teacher's trajectory: the numerical ODE path from
$z_0$ to $x_0$ is generally curved, and this is the straight chord between its two ends.

\begin{table}[h]
\centering
\caption{Coupling, at one target and at four. Three seeds each. Both rows of a pair share the same
caption-to-image map, so that map cancels in the difference.}
\label{tab:coupling}
\small
\begin{tabular}{@{}llrrl@{}}
\toprule
Targets per caption & Noise & CompBench++ & GenEval2 & Gain from coupling \\
\midrule
one, fixed & fresh & \obs{0.2696} & \obs{0.1224} & --- \\
one, fixed & \textbf{its own} & \obs{\textbf{0.3528}} & \obs{0.1610} & \obs{$+0.0832$ $[+0.0608,+0.1049]$} \\
four, maximum churn & fresh & \obs{0.2429} & \obs{0.0857} & --- \\
four, maximum churn & \textbf{each its own} & \obs{\textbf{0.3293}} & \obs{0.1338} & \obs{$+0.0863$ $[+0.0546,+0.1194]$} \\
\bottomrule
\end{tabular}
\end{table}

\paragraph{The one-line result.}
Keeping \textbf{all four} teacher images, each paired with its own noise, scores \obs{0.3293}.
Keeping \textbf{one} image with random noise scores \obs{0.2696}. \textbf{The problem the first
two stages were circling --- why using fewer teacher images helped --- was never that there were
several images. It was that the loss had thrown away which noise made which image.}

\paragraph{Coupling is worth more than choosing, ordering or fixing the target.}
It is \obs{$+0.083$}, against \obs{$+0.042$} for fixing the target at this budget and
\obs{$+0.034$} for the VQAScore top-1. It is also the best identified of these numbers, because the two
arms of each pair share one caption-to-image map (\cref{sec:assign}).

\paragraph{Coupling and the number of targets are close to additive.}
The gain from coupling is almost the same at one target (\obs{$+0.0832$}) and at four
(\obs{$+0.0863$}); the interaction is \obs{$+0.0032$}. Equivalently, the cost of using four targets
rather than one is \obs{$-0.0266$} $[-0.0581,+0.0026]$ without coupling and \obs{$-0.0235$}
$[-0.0464,+0.0008]$ with it. \textbf{Coupling does not explain away the benefit of using fewer
targets, and it does not remove it.} Four coupled targets still score below one coupled target.

An earlier version of this report claimed all four targets could be kept ``at no cost in
alignment''. That is contradicted by \obs{0.3293} against \obs{0.3528}. The supported statement is:
correct coupling gives a large gain that is nearly independent of how many targets are used, and
fewer targets separately gives a smaller one. Whether the \obs{$-0.0235$} is worth paying depends on
the diversity measurement, which is not yet done (\cref{sec:limits}).

\paragraph{This is not a new algorithm.}
Pairing a generated endpoint with the noise that produced it, and training on the straight
interpolation between them, is \textbf{Reflow}, introduced with rectified flow and used at scale by
InstaFlow. Related coupling constructions appear in Multisample Flow Matching and StraightFM. Our
contribution here is not the method: it is the measurement that this coupling, and not target
stability or label quality, accounts for most of what we had attributed to selection, together with
the controls in \cref{sec:stage2} that rule out the alternative. Read this report as an empirical
diagnosis of a known mechanism in a setting where it had been discarded by accident.

\paragraph{What is specific to our setting.}
Reflow normally couples a model's own samples. Here the endpoints are \emph{teacher} samples that
have been \emph{selected} by a reward, and the same caption is revisited many times. Selecting
best-of-$N$ by a reward biases which source noises survive, so the source marginal seen by the
student is no longer the Gaussian that Reflow assumes. We have not measured that bias, and it is
the part of this problem we think is actually open.

\paragraph{Predicted before it was measured.}
We split the flow-matching gradient variance by the law of total variance over the four candidates
$j$, on \obs{100} captions, using a fixed random \obs{200{,}000}-coordinate sketch:
\begin{equation}
  \operatorname{Var}(g)=\underbrace{\mathbb E_j\!\left[\operatorname{Var}(g\mid j)\right]}_{\text{noise inside one target}}
  +\underbrace{\operatorname{Var}_j\!\left(\mathbb E[g\mid j]\right)}_{\text{switching between targets}} .
\label{eq:varsplit}
\end{equation}

\begin{center}
\small
\begin{tabular}{@{}lrrr@{}}
\toprule
Noise & within a target & switching & switching share \\
\midrule
fresh & \obs{6.00e-4} & \obs{7.59e-5} & \obs{11.2\%} \\
coupled & \obs{4.16e-4} & \obs{8.36e-5} & \obs{16.7\%} \\
\bottomrule
\end{tabular}
\end{center}

Switching accounts for only \obs{11.2\%} of the variance, so removing it cannot explain a large
effect. Coupling cuts the \emph{dominant} term by \obs{31\%}, and the total by \obs{26\%} --- about
\obs{2.4} times what perfect persistence removes. We recorded the prediction that coupling would
therefore beat persistence before those runs finished. Measured: \obs{$+0.083$} against
\obs{$+0.042$}, a ratio of \obs{2.0}.


\section{Scope: every effect here depends on budget and on how often a caption repeats}
\label{sec:scope}

The effects above were measured at \obs{6{,}000} updates over \obs{3{,}000} captions, which is
\obs{8} visits per caption. Neither number is neutral. Changing the training budget and changing the
number of unique captions each move the results, and they move them in \emph{opposite} directions.

\subsection{Persistence needs revisits, and disappears without them}

A fixed target can only differ from a redrawn one on a caption's second visit. We held the number of
updates fixed and changed only how many unique captions they are spread over.

\begin{table}[h]
\centering
\caption{Fixed against redrawn target, at three pool sizes and one budget
(\obs{6{,}000} updates $\times$ \obs{4} workers = \obs{24{,}000} caption visits).}
\label{tab:captioncount}
\small
\begin{tabular}{@{}lrrrl@{}}
\toprule
Captions & Visits each & Fixed & Redrawn & Gap \\
\midrule
\obs{3{,}000}  & \obs{8.0} & \obs{0.2679} & \obs{0.2376} & \obs{$+0.0303$} $[+0.0025,+0.0574]$ \\
\obs{12{,}000} & \obs{2.0} & \obs{0.2466} & \obs{0.2355} & \obs{$+0.0111$} $[+0.0001,+0.0209]$ \\
\obs{23{,}999} & \obs{1.0} & \obs{0.2450} & \obs{0.2438} & \bad{$+0.0012$} $[-0.0153,+0.0174]$ \\
\bottomrule
\end{tabular}
\end{table}

The gap falls monotonically and reaches zero exactly where it must: at one visit per caption the two
arms are the \emph{same procedure}, so \obs{$+0.0012$} is this experiment's own noise floor and it
validates the design. \textbf{The persistence result belongs to low-data, repeatedly recycled
distillation.} It is not a general principle of few-step distillation, and a paper claiming it as one
would be wrong.

\subsection{Coupling, unlike persistence, does not shrink with pool size}

Persistence decays across \cref{tab:captioncount} because its mechanism --- a target that stops
changing --- needs revisits. Coupling has no such requirement, and it does not decay.

\begin{table}[h]
\centering
\caption{Own source noise against fresh, at two pool sizes and one budget
(\obs{6{,}000} updates). Reference-only: see the map caveat in \cref{sec:assign}.}
\label{tab:couplingscope}
\small
\begin{tabular}{@{}lrrrl@{}}
\toprule
Captions & Visits each & Fresh & Own source & Gain \\
\midrule
\obs{3{,}000}  & \obs{8.0} & \obs{0.2679} ($n{=}3$) & \obs{0.3594} ($n{=}3$) & \obs{$+0.0915$} \\
\obs{12{,}000} & \obs{2.0} & \obs{0.2466} ($n{=}2$) & \obs{0.3289} ($n{=}3$) & \obs{$+0.0823$} \\
\obs{23{,}999} & \obs{1.0} & \obs{0.2450} ($n{=}2$) & \obs{0.3267} ($n{=}3$) & \obs{$+0.0817$} \\
\bottomrule
\end{tabular}
\end{table}

The gain is nearly flat across an \obs{8}-fold change in visits per caption, and at \obs{1.0} visits
persistence has no mechanism at all --- yet \obs{$+0.0817$} of it is still present. \textbf{So the
drop from \obs{$+0.0832$} at \obs{6}k updates to \obs{$+0.0310$} at \obs{40}k is a \emph{budget}
effect, not a data-scope effect, and it is not mainly persistence.} An earlier draft of this report
described it as a data-scope effect; that was wrong.

The residual gradient \obs{$+0.0915\to+0.0823\to+0.0817$} points the right way for a small
persistence component --- \obs{3{,}000} captions is the only row where a caption is revisited --- and
would put that component near \obs{$+0.010$} on top of a visit-independent \obs{$+0.082$}. We do not
claim it: two of the three fresh cells are $n{=}2$, and \obs{0.010} is not separable from run noise
here. It is a hypothesis for the controlled ladder, not a measurement.

Two cautions. These two cells are the ones whose fresh and coupled arms used different
caption-to-image maps (\cref{sec:assign}), so they are reference-only and their intervals are wider
than a paired design would give. And a flat gain across pool size is here inseparable from a flat
gain across revisit rate, because the two move together by construction: breaking that would need a
cell at, say, \obs{3{,}000} captions and \obs{750} updates.

\subsection{Coupling shrinks with budget; the label grows}

\begin{table}[h]
\centering
\caption{The same three factors at two budgets. They do not move together.}
\label{tab:budgetscope}
\small
\begin{tabular}{@{}lrr@{}}
\toprule
Factor & \obs{6{,}000} steps & \obs{40{,}000} steps \\
\midrule
Label, $\rho\;\obs{0}\to\obs{1}$, target fixed & \bad{$+0.0042$} (not detected) & \obs{$+0.0366$} \\
Persistence, target stops changing & \obs{$+0.0417$} & \obs{$+0.0954$} \\
\textbf{Coupling}, one target & \obs{$\mathbf{+0.0832}$} & \obs{$\mathbf{+0.0310}$} $[+0.0178,+0.0432]$ \\
\bottomrule
\end{tabular}
\end{table}

Coupling is the only factor that gets \emph{smaller} with more training, and it more than halves.
At \obs{40{,}000} steps it also stops being the best option: a coupled single target reaches
\obs{0.4074}, against \obs{0.4101} for a PRS label and \obs{0.4131} for a VQAScore label on the same
budget. The difference between coupling and a PRS label there is \obs{$-0.0026$}
$[-0.0149,+0.0092]$ --- nothing.

\paragraph{What we think this means.}
Coupling gives the student an easier regression problem: the interpolation it must invert is the
straight chord of a pair the teacher realised, rather than a chord between unrelated points. An
easier problem helps most when training is short and matters less once the student has had time to
fit the harder one. On this reading coupling is an \emph{optimisation} aid rather than a change in
what is learned, which is consistent with it shrinking while the label effect grows.

\paragraph{So the headline must be stated with its scope.}
Coupling is a large win at a short budget, worth about \obs{$+0.083$} at \obs{6{,}000} steps. By
\obs{40{,}000} steps it is worth \obs{$+0.031$} and is matched by simply using a label. We do not
have evidence that it helps at budgets longer than that.

\subsection{Most of the coupling gain is that the noise stops changing, not that it is the right noise}

If coupling works because the source and the endpoint belong together, then pairing an endpoint with
some \emph{other} fixed noise should not help. We tested that directly, holding the target fixed and
changing only which noise it is paired with.

\begin{table}[h]
\centering
\caption{Same fixed target; only the paired noise differs. \obs{6{,}000} steps.}
\label{tab:noiseid}
\small
\begin{tabular}{@{}llrl@{}}
\toprule
Noise & What it is & CompBench++ & Step \\
\midrule
fresh & a new draw on every visit & \obs{0.2696} & --- \\
another caption's & a fixed draw, unrelated to this endpoint & \obs{0.3212}$^{\dagger}$ & \obs{$+0.0516$} $[+0.0333,+0.0672]$ \\
its own & the source the teacher actually used & \obs{0.3528} & \obs{$+0.0316$} $[+0.0172,+0.0488]$ \\
\bottomrule
\end{tabular}

\vspace{2pt}
{\footnotesize $^{\dagger}$one seed; the remaining seeds and two further controls
(another candidate of the same caption, and a fixed draw that was never a rollout source) are still
running.}
\end{table}

The split is roughly \obs{60}/\obs{40}. Most of the benefit comes from the noise simply being
\emph{fixed per target}: an unrelated noise from another caption recovers \obs{$+0.052$} of the
\obs{$+0.083$}. Being the \emph{correct} source is worth a further \obs{$+0.032$} --- real, and
comparable to the whole label effect at \obs{40{,}000} steps, but not the main part.

\textbf{This narrows the claim.} An earlier version of this report said the loss had ``thrown away
which noise made which image''. That is only \obs{40\%} of the story. The larger part is that the
loss redrew the source on every visit, and stopping that redraw is most of the gain.

\section{Diversity: the two ways of gaining alignment are not equivalent}
\label{sec:diversity}

Every alignment number here uses one image per prompt, so none of them can see a model that has
learned to draw the same safe picture every time. A pre-registered gate from an earlier campaign in
this project measures that directly: \obs{4} images per prompt differing only in initial noise, mean
pairwise DINOv2 cosine distance and mean pairwise LPIPS, failing if the paired interval's upper
bound falls below \obs{$-0.02$}. We ran the same harness, unchanged, on our arms.

\begin{table}[h]
\centering
\caption{Conditional diversity: how different are \obs{4} images of the same caption? Higher is more
varied. \obs{500} prompts.}
\label{tab:diversity}
\small
\begin{tabular}{@{}llrr@{}}
\toprule
Arm & Budget & DINO $1-\cos$ & LPIPS \\
\midrule
untrained & --- & \obs{0.5942} & \obs{0.4905} \\
\midrule
\texttt{RAND-soft}, target moves & \obs{40}k & \obs{0.3909} & \obs{0.3339} \\
\texttt{RAND-fix}, target fixed & \obs{40}k & \bad{0.2236} & \bad{0.2593} \\
\texttt{PRS-top1} & \obs{40}k & \bad{0.2004} & \obs{0.2720} \\
\texttt{VQA-top1} & \obs{40}k & \bad{0.1994} & \obs{0.2670} \\
\midrule
fresh noise, one target & \obs{6}k & \obs{0.4069} & \obs{0.3370} \\
coupled, one target & \obs{6}k & \obs{0.4609} & \obs{0.6087} \\
fresh noise, four targets & \obs{6}k & \obs{0.4580} & \obs{0.3720} \\
\textbf{coupled, four targets} & \obs{6}k & \obs{\textbf{0.4618}} & \obs{\textbf{0.5871}} \\
\bottomrule
\end{tabular}
\end{table}

\paragraph{Fixing the target fails the gate.}
\texttt{RAND-fix} $-$ \texttt{RAND-soft} is \bad{$-0.1673$} $[-0.1761,-0.1589]$ on DINO and
\bad{$-0.0746$} $[-0.0795,-0.0698]$ on LPIPS, against a failure threshold of \obs{$-0.02$}. The
\obs{$+0.0954$} of alignment that a fixed target buys is paid for by losing \obs{43\%} of the
variation the model had left. Adding a label collapses it further.

\paragraph{Coupling passes it.}
Coupled four targets $-$ fresh one target is \obs{$+0.0549$} $[+0.0447,+0.0649]$ on DINO and
\obs{$+0.2501$} $[+0.2431,+0.2572]$ on LPIPS. Coupling raises alignment \emph{and} variation.

We checked the images before believing the LPIPS figure, because a value above the untrained model's
is what colour corruption would also produce. It is not corruption: the coupled columns are sharper
and carry more detail, clipped pixels stay at \obs{3}--\obs{8\%} against \obs{55\%} for the failed
student of \cref{sec:failures}, and high-frequency energy rises only slightly.

\paragraph{This is the sharpest distinction in the report.}
An earlier campaign here measured alignment against DINO diversity across twelve arms and found them
trading at $r=\obs{-0.9622}$. \texttt{RAND-fix} obeys that line exactly. Coupling does not: it moves
off it. So the two routes to alignment differ in kind, not only in size --- one narrows what a
caption can produce and the other does not.

\paragraph{Limit.}
The coupled arms here are \obs{6{,}000}-step models and the failing arms are \obs{40{,}000}-step
models. Less training leaves more variation, so the comparison across those two blocks is
indicative; the \obs{2}$\times$\obs{2} inside the \obs{6}k block is matched and is what the claim
rests on. Coupled arms at \obs{40{,}000} steps are training, and \cref{sec:scope} already shows the
alignment advantage shrinks there.

\section{The five scores}
\label{sec:scores}

Four of these compare a generated image with \textbf{the real photograph the caption was written
for}. COCO gives us that photograph, so a generated image can be scored by how close it is to the
one a person actually described. The fifth compares the image with the caption directly.

\begin{table}[h]
\centering
\caption{What each score is and what it costs. ``Decode?'' is whether the image must be turned from
a latent into pixels before the score can be computed.}
\label{tab:scoredefs}
\small
\begin{tabularx}{\textwidth}{@{}llX@{}}
\toprule
Score & Decode? & What it computes \\
\midrule
\textbf{VQAScore} & yes & Decode the image, then ask an 11-billion-parameter vision--language model
  (\texttt{clip-flant5-xxl}) how well the image matches the \emph{caption}. The only score here that
  reads the caption. We treat it as ground truth for ``which image is best''. \\
\textbf{CLIP-I} & yes & Decode the image, embed it with the CLIP image encoder
  (\texttt{ViT-L/14}), embed the real photograph the same way, take the cosine between them. The
  ``-I'' marks it as \emph{image-to-image}: unlike CLIPScore, it never looks at the caption. \\
\textbf{DINOv2} & yes & The same idea with DINOv2 features (\texttt{dinov2-base}, CLS token) instead
  of CLIP. DINOv2 is trained without text, so it responds to layout and object identity rather than
  to anything caption-like. \\
\textbf{LPIPS} & yes & A perceptual \emph{distance} between the decoded image and the photograph,
  from AlexNet features. We negate it so that, like the others, larger means better. \\
\textbf{PRS} (ours) & \textbf{no} & Paired-reference similarity. Take the denoiser's own hidden
  states, averaged over image tokens, for the generated latent and for the encoded photograph at one
  noise level, and take the cosine. It needs no decoder and no second network, because the denoiser
  has already computed these states. \\
\bottomrule
\end{tabularx}
\end{table}

Formally, with $\overline{h}$ the token average of one denoiser block, $\mathcal E$ the VAE encoder
and $x_\sigma$ the noising of \cref{eq:path},
\begin{equation}
  \mathrm{PRS}_j = \cos\Bigl(\overline{h}\bigl(x_\sigma(z_j),\,T\sigma,\,c\bigr),\;
  \overline{h}\bigl(x_\sigma(\mathcal E(x_{\text{ref}})),\,T\sigma,\,c\bigr)\Bigr).
\label{eq:prs}
\end{equation}

\paragraph{How well each one ranks.}
We score four candidates per caption on \obs{3{,}000} captions and ask what share of the available
gain an argmax over that score recovers. Writing $R_j$ for the VQAScore of candidate $j$ and
$\hat\jmath$ for the one a score picks,
\begin{equation}
  \rho=\frac{\sum_{\text{captions}}\bigl(R_{\hat\jmath}-\tfrac1N\sum_j R_j\bigr)}
            {\sum_{\text{captions}}\bigl(\max_j R_j-\tfrac1N\sum_j R_j\bigr)} ,
\label{eq:rho}
\end{equation}
a ratio of sums, so each caption carries its own available gain as weight. A perfect ranker gives
$\rho=1$, a random one $\rho=0$.

\begin{table}[h]
\centering
\caption{Left: ranking ability on one shared pool, so all five see the same images and are graded
against the same VQAScore table. Right: what happens when that score is used to pick the training
target, all arms on the same candidates, three seeds each.}
\label{tab:rankers}
\small
\begin{tabular}{@{}lrlrl@{}}
\toprule
Score & $\rho$ & PRS $-$ this, ranking & CompBench++ & this $-$ no score, trained \\
\midrule
CLIP-I & \obs{40.64\%} & \bad{$+1.20$ $[-4.16,+6.55]$} & \obs{\textbf{0.3094}} & \obs{$+0.0249$ $[+0.0115,+0.0368]$} \\
\textbf{PRS} & \obs{\textbf{41.84\%}} & --- & \obs{0.2876} & \bad{$+0.0031$ $[-0.0202,+0.0235]$} \\
DINOv2 & \obs{39.08\%} & \bad{$+2.77$ $[-2.46,+8.02]$} & \obs{0.2869} & \bad{$+0.0024$ $[-0.0162,+0.0214]$} \\
LPIPS & \obs{23.26\%} & \obs{$+18.58$ $[+12.85,+24.38]$} & --- & --- \\
no score & \obs{0\%} & --- & \obs{0.2845} & --- \\
\bottomrule
\end{tabular}
\end{table}

\textbf{Only CLIP-I is worth using.} PRS and DINOv2 cannot be told apart from using no score at all,
and PRS is worse than CLIP-I by \obs{$-0.0218$} $[-0.0428,-0.0056]$.

\paragraph{Ranking ability does not predict usefulness.}
PRS has the best $\rho$ of the three and produces the second-worst model. The offline measure and the
trained outcome disagree in \emph{order}, not merely in size, so $\rho$ should not be used to choose
a score.

\paragraph{The tuning did not fully replicate.}
PRS has two settings, the noise level and which denoiser block to read. Sweeping \obs{11}$\times$
\obs{9} of them on a separate pool of \obs{500} captions gave a best cell of \obs{45.84\%}; the same
setting on the disjoint \obs{3{,}000}-caption pool gives \obs{41.84\%}. A drop of \obs{4} points is
what choosing the best of \obs{99} cells on \obs{500} captions looks like.

\paragraph{And ``no decode'' does not make it cheap.}
The decoder PRS skips costs \obs{60.8}\,ms; the two extra denoiser calls it adds cost
\obs{83.0}\,ms. \textbf{The step we skip is cheaper than the step we add.}

\begin{table}[h]
\centering
\caption{Cost per caption in milliseconds: four candidates, \obs{8} teacher steps, one H100, mean of
\obs{19} captions after a warm-up. ``Shared'' is text encoding plus the four teacher rollouts,
\obs{847.9}\,ms, which every score pays before any of them starts.}
\label{tab:cost}
\small
\begin{tabular}{@{}lrrrr@{}}
\toprule
Score & Needs a decode? & Score only & Shared $+$ score & GPU-hours / 1000 captions \\
\midrule
LPIPS & yes & \obs{73.2} & \obs{921.1} & \obs{0.2559} \\
PRS & no & \obs{94.7} & \obs{942.5} & \obs{0.2618} \\
DINOv2 & yes & \obs{104.1} & \obs{952.0} & \obs{0.2644} \\
CLIP-I & yes & \obs{120.8} & \obs{968.6} & \obs{0.2691} \\
VQAScore (11B) & yes & \obs{340.8} & \obs{1188.7} & \obs{0.3302} \\
\bottomrule
\end{tabular}
\end{table}

Against VQAScore, PRS is \obs{3.6} times cheaper on the score alone but only \bad{\obs{1.26}} times
cheaper end to end, because every score must first make the four images. Against LPIPS it is more
expensive outright.

\section{The training recipe, and what it trades}
\label{sec:objective}

With the label held identical, moving from the flow-matching recipe to the consistency-distillation
recipe is worth \obs{$+0.2043$} $[+0.1770,+0.2254]$ under the VQAScore label and \obs{$+0.1879$}
$[+0.1715,+0.2071]$ under PRS. That is the largest number in this report.

\paragraph{This is a recipe difference, not an ``objective'' effect.}
The two groups differ in more than the loss. They use different training states (endpoint versus the
teacher's trajectory), a different time support, a different source--target construction, a different
prediction target, and a different distance (MSE against pseudo-Huber). We cannot attribute
\obs{$+0.20$} to any one of those. Factorising them is the obvious next study and we have not run
it.

It is not a free win. GenEval2 stores a score for every atom of every prompt, so the two poolings and
the skill breakdown can be recomputed from what is on disk. The geometric mean asks ``did it get
everything right?''; the arithmetic mean asks ``how much did it get right?''.

\begin{table}[h]
\centering
\caption{GenEval2 re-pooled from stored atom scores. The geometric mean reproduces the recorded
value for \obs{71/71} results, which validates the pooling.}
\label{tab:skills}
\small
\begin{tabular}{@{}lrrrrrrr@{}}
\toprule
Group & GM & AM & object & attribute & position & count & verb \\
\midrule
CD, \obs{6}k & \obs{0.1912} & \obs{0.5548} & \obs{\textbf{0.7526}} & \obs{0.6262} & \obs{0.4216} & \obs{0.3317} & \obs{0.1997} \\
FM, \obs{40}k, fixed & \obs{0.2196} & \obs{0.5795} & \obs{0.7201} & \obs{\textbf{0.6496}} & \obs{\textbf{0.5202}} & \obs{\textbf{0.3696}} & \obs{\textbf{0.2515}} \\
FM, \obs{6}k, fixed & \obs{0.1318} & \obs{0.4566} & \obs{0.5467} & \obs{0.5540} & \obs{0.4190} & \obs{0.2712} & \obs{0.1040} \\
\bottomrule
\end{tabular}
\end{table}

The consistency objective draws \emph{objects} better; flow matching at the longer budget is better
at every \emph{relational} skill, and the gap on \texttt{position} is \obs{0.099}. T2I-CompBench++ is
dominated by attribute binding on well-drawn objects, which is the consistency objective's strength,
so the headline understates what flow matching is good at. Partial credit is large everywhere:
AM $\approx\obs{0.58}$ against GM $\approx\obs{0.22}$.


\section{Absorbing guidance on its own}
\label{sec:g8}

Our students learn two things at once: one network call per step instead of two, and \obs{4} steps
instead of \obs{8}. Trying to keep guidance \emph{inside} a 4-step student failed three times
(\cref{sec:failures}). Doing the first job alone works.

We trained $G$ to reproduce the guided teacher with a single conditional call at the same \obs{8}
steps. The target is free: Euler's own update gives the guided speed the teacher used, and we record
it in float32 as the teacher computes it. Rebuilding it later by subtracting stored states would
lose several percent, because those states are kept in bfloat16 and the difference between two
neighbours is much smaller than either.

\begin{center}
\small
\begin{tabular}{@{}lrrr@{}}
\toprule
 & untrained, 1 call & $G$, 1 call & teacher, 2 calls \\
\midrule
relative $L_2$ to the teacher & \obs{0.7730} & \obs{0.3339} & \obs{0} \\
cosine to the teacher & \obs{0.6258} & \obs{0.9376} & \obs{1} \\
VQAScore & \obs{0.7155} & \obs{0.8362} & \obs{0.8459} \\
\bottomrule
\end{tabular}
\end{center}

$G$ removes \obs{56.8\%} of the speed error, keeps \obs{93\%} of the teacher's VQAScore advantage
over the unguided model, and on the benchmarks matches it outright (\cref{tab:refs}) --- at half the
network calls.


\section{Uncertainty}
\label{sec:assign}

\paragraph{The caption-to-image map is a variance component, and we missed it at first.}
A \texttt{fix} arm draws one candidate per caption from a seed that depends on the caption only. So
all three seeds of an arm train on the \emph{same} map, and a bootstrap over seeds never resamples
it. Running three more \texttt{RAND-fix} arms with different maps as well as different seeds:

\begin{center}
\small
\begin{tabular}{@{}lrr@{}}
\toprule
 & mean & standard deviation \\
\midrule
one map, three seeds & \obs{0.2857} & \obs{0.0115} \\
three maps, three seeds & \obs{0.2770} & \obs{0.0212} \\
\bottomrule
\end{tabular}
\end{center}

The map contributes about \obs{0.018}, more than the optimiser noise it hid behind, so intervals on
fixed-target arms were roughly \obs{1.8} times too narrow. Corrected, the \obs{6}k persistence effect
is about \obs{$+0.042$} $[+0.012,+0.072]$ --- still clear of zero, and weaker than first reported.

\textbf{The coupling result is better protected}: both arms of each pair in \cref{tab:coupling} use
the same map --- verified from each run's resolved configuration, not from its name --- so the
map's main effect cancels in the difference. That does not rule out a map-by-coupling interaction,
which our design cannot see. Repeating the coupling comparison across several independent maps is
the check we owe it.

\paragraph{Two pairing defects found afterwards, with different scopes.}
\emph{The caption-to-image map.} A fresh arm launched \emph{without} \texttt{-{}-target\_schedule}
selects its target with \texttt{assign\_rng(seed,idx).choice}, while an arm launched with
\texttt{one\_fixed, m{=}1} uses \texttt{assign\_rng(seed,idx).permutation(N)[:1]}. Same generator
stream, different call: they agree on \obs{5991}/\obs{23999} captions, or \obs{0.250}, which is
chance for $N=4$. This does \emph{not} affect \cref{tab:coupling}, whose arms both passed the
schedule. It does affect the scope cells of \cref{sec:scope}, whose fresh arms did not --- so those
two cells carry an extra $\sqrt2\times\obs{0.018}\approx\obs{0.026}$ of map uncertainty and are
reference-only rather than matched comparisons.

\emph{The noise levels.} The flow-matching $\sigma$ draw comes from the global generator, and
\texttt{-{}-coupling fresh} then draws its noise from that same generator while every coupled mode
uses a private one. A fresh arm therefore consumes more of the global stream per update, and the
two arms walk different $\sigma$ sequences from update \obs{2} onward. This follows from the
coupling mode rather than the schedule, so unlike the map defect it touches \emph{every}
fresh-versus-coupled contrast here, \cref{tab:coupling} included. It does not bias any estimate ---
$\sigma$ is drawn identically in distribution either way --- but it is unpaired noise in exactly the
contrast being measured. Both are fixed for future runs (\texttt{-{}-paired\_sigma}, and a target
map that is materialised and hash-checked before launch rather than assumed from the flags).

\paragraph{No null here is a claim of absence.}
An interval crossing zero means we failed to detect an effect. To claim ``no effect'' at a margin of
\obs{$\pm0.01$} would need roughly \obs{44} seeds per arm; we ran three. The order nulls in
\cref{sec:stage2} are strong because churn was varied sevenfold and both intervals are tight around
zero, but they are still failures to detect.


\section{What failed}
\label{sec:failures}

\paragraph{Keeping guidance inside a 4-step student.}
Three attempts, all worse than not training at all (\obs{0.2557}): supervising only the guided
mixture gave \bad{0.14--0.16}; giving each branch its own target through a common shift gave
\bad{0.1165}; sampling the states from the student's own path gave \bad{0.1889}. A diagnostic ruled
out loss scale, gradient magnitude and gradient cancellation: both losses start at the same value
(\obs{78.87} against \obs{78.91}), the branch gradient is \obs{6.7} times \emph{smaller}, and the two
branch gradients agree rather than cancel (cosine \obs{$+0.606$}). The last attempt produces correct
shapes in destroyed colour: its latent is no larger than the untrained model's
(\obs{5.16} against \obs{6.28}) but its channels collapse together --- spread of channel means
\obs{1.348}$\to$\obs{0.675} --- so \obs{40\%} of the decoder output leaves $[-1,1]$. Its score comes
almost entirely from the caption-model categories, while every detector category is near zero and
spatial is exactly \obs{0.0000}. We use the two-stage separation instead (\cref{sec:g8}).

\paragraph{A 28-step teacher.}
Better images make selection pointless: $\rho$ for PRS falls from \obs{37.21\%} to \obs{1.17\%}, the
gain available per caption from \obs{0.0893} to \obs{0.0416}, and ties rise from \obs{20.4\%} to
\obs{44.3\%}. \textbf{Selection only helps when the candidates differ.}


\section{Limits}
\label{sec:limits}

\begin{itemize}
  \item \textbf{Same seed does not mean same image across hardware.} Rebuilding a cache with an
    identical pool, seeds, steps and guidance reproduced the random index for \obs{3000/3000}
    captions, but VQAScore moved by up to \obs{1.3e-2} and the best candidate changed for
    \obs{116/3000} (\obs{3.9\%}). Runs landed on different GPU models and bfloat16 reductions differ.
    Any comparison spanning two cache builds carries about \obs{4\%} label noise, which is why every
    ranker in \cref{tab:rankers} is graded against one table.
  \item \textbf{Score and appearance can move apart.} The failed student in \cref{sec:failures}
    scored higher than a worse-looking one, because caption models see through colour damage that
    detectors do not. Any surprising score deserves a look at the images.
  \item \textbf{One image per prompt} at test time, so each number carries sampling noise of its own.
  \item \textbf{Relations are not learned.} Every arm improves how objects are drawn. None fixes
    spatial or counting relations, which is where the largest headroom remains.
  \item \textbf{Diversity is not yet measured.} A fixed or coupled target could improve one-sample
    alignment while narrowing the range of images a caption can produce. Generation for that
    measurement is running and is not reported here.
\end{itemize}


\section{What to use}

\begin{enumerate}
  \item \textbf{Select at sampling time, not in the weights.} Best of \obs{4} by VQAScore adds
    \obs{$+0.038$} with no training, against \obs{$+0.034$} from the same signal as a label after
    \obs{40{,}000} updates. It also lifts \texttt{spatial}, which no trained arm moved.
  \item \textbf{Pair every target with its own noise} if the training budget is short. Worth
    \obs{$+0.083$} at \obs{6{,}000} steps, \obs{$+0.031$} at \obs{40{,}000}, and nothing over a plain
    label by then. It costs nothing to implement.
  \item \textbf{Do not fix the target to buy alignment.} It works, and it fails the diversity gate at
    \bad{$-0.167$} DINO. Coupling reaches similar alignment without that cost.
  \item \textbf{Do not spend on the ranker.} Only CLIP-I beat using no score at all. PRS and DINOv2
    did not.
  \item \textbf{Absorb guidance separately} from step reduction: half the calls, no loss
    (\cref{sec:g8}).
  \item \textbf{Measure against $G$ at \obs{4} steps} (\obs{0.2859}), and report the \obs{28}-step
    generator (\obs{0.5053}), not an untrained model.
\end{enumerate}

\paragraph{What we would not repeat.}
Measuring at one budget and one pool size. Three of the conclusions in earlier drafts of this report
reversed or vanished when the same arms were run at \obs{40{,}000} steps or on \obs{24{,}000}
captions. Treating $\rho$ as a forecast. Drawing the caption-to-image map from a seed that does not
vary with the run. And editing the trainer while a batch is queued, which splits the batch silently
across two versions of the code.

\section{Files}

\begin{table}[h]
\centering
\small
\begin{tabular}{@{}ll@{}}
\toprule
Trainer, both objectives, coupling, target schedules & \texttt{phaseC/train\_pilot.py} \\
Candidate cache, PRS and the ranker baselines & \texttt{phaseN/build\_coco\_selection.py} \\
Guidance distillation and its gate & \texttt{phaseN/train\_guidance\_distill.py}, \texttt{verify\_g8.py} \\
Gradient variance split & \texttt{phaseN/grad\_variance.py} \\
Results and paired bootstrap & \texttt{phaseN/analyze\_arms.py}, \texttt{report\_experiments.py} \\
GenEval2 re-pooling and skills & \texttt{phaseN/geneval2\_repool.py} \\
Label timings & \texttt{phaseN/label\_cost.py} \\
\bottomrule
\end{tabular}
\end{table}


\end{document}
